\documentclass[11pt]{article}

\usepackage[margin=1in]{geometry}
\usepackage[T1]{fontenc}
\usepackage[utf8]{inputenc}
\usepackage{amsmath,amssymb,bm}
\usepackage{booktabs}
\usepackage{array}
\usepackage{longtable}
\usepackage{xcolor}
\usepackage{hyperref}

\setlength{\tabcolsep}{4pt}

\hypersetup{
    colorlinks=true,
    linkcolor=blue,
    citecolor=blue,
    urlcolor=blue
}

\newcommand{\ket}[1]{\left|#1\right\rangle}
\newcommand{\bra}[1]{\left\langle #1\right|}
\newcommand{\op}[2]{\left|#1\right\rangle\left\langle #2\right|}
\newcommand{\needcite}{\textbf{[need citation]}}
\newcommand{\testbox}[1]{%
    \par\medskip
    \noindent
    \fcolorbox{black}{gray!8}{%
        \begin{minipage}{0.96\linewidth}
        \textbf{Test point.} #1
        \end{minipage}
    }
    \par\medskip
}
\newcommand{\decisionbox}[1]{%
    \par\medskip
    \noindent
    \fcolorbox{black}{yellow!12}{%
        \begin{minipage}{0.96\linewidth}
        \textbf{Decision criterion.} #1
        \end{minipage}
    }
    \par\medskip
}

\title{Scientific Validation Plan for the Bright--Dark Model Used to Probe the Spinorial Sector of a Spin-Peierls Chain}
\author{}
\date{\today}

\begin{document}

\maketitle

\section{Purpose of the plan}

The goal is not only to produce spectra. The goal is to build a sequence of tests that can answer, in a falsifiable way, the central question:
\begin{quote}
Can multidimensional optical spectroscopy, through an effective spin-orbit coupling, probe the dynamics of the spin sector in a Spin-Peierls chain?
\end{quote}

The current model is an effective model. It does not yet contain a complete microscopic Hamiltonian with \(d\) orbitals, crystal field, explicit spin-orbit coupling, phonons, and a dynamic spin chain. The strategy must therefore separate three levels:
\begin{enumerate}
    \item \textbf{What the solver actually computes}: a multidimensional optical response of a bright--dark--mode Hamiltonian.
    \item \textbf{What the parameters encode}: \(\delta\), \(C_1\), \(V_{\mathrm{BD}}\), \(\omega_Q(k)\), \(g_Q(k)\).
    \item \textbf{What we want to infer physically}: spin correlations, Spin-Peierls dynamics, magnetic-field effects, bright--dark coherence, and the role of SOC.
\end{enumerate}

The strategy is to add one physical layer at a time. Each layer must have a control in which the expected signature disappears. If a signature is already present in a simpler model, it cannot be assigned unambiguously to the more complex mechanism.

\section{Reference model to validate}

The target effective Hamiltonian is
\begin{equation}
H(k)
=
E_D\op{D}{D}
+
E_B\op{B}{B}
+
\omega_Q(k)a^\dagger a
+
V_{\mathrm{BD}}^0 L_{\mathrm{BD}}
+
g_Q(k)L_{\mathrm{BD}}(a+a^\dagger),
\end{equation}
with
\begin{equation}
L_{\mathrm{BD}}
=
\op{D}{B}+\op{B}{D},
\end{equation}
and
\begin{equation}
V_{\mathrm{BD}}^0
=
V_0+\lambda_\delta\delta+\lambda_C C_1.
\end{equation}

The dispersive part is
\begin{equation}
\omega_Q(k)
=
\omega_Q
+
s\left[\omega_{\mathrm{SP}}(k)-\omega_{\mathrm{SP}}(0)\right],
\end{equation}
and the dynamic coupling is
\begin{equation}
g_Q(k)
=
g_Q\left(1+A_Q\cos k\right).
\end{equation}

The model should be interpreted as an inference chain:
\begin{equation}
\text{spin chain}
\rightarrow
(\delta,C_1,\Delta_{\mathrm{spin}})
\rightarrow
(V_{\mathrm{BD}},\omega_Q(k),g_Q(k))
\rightarrow
S(\omega_1,\omega_3).
\end{equation}

\section{Scientific questions and target observables}

\begin{longtable}{>{\centering\arraybackslash}p{0.07\linewidth} p{0.34\linewidth} p{0.27\linewidth} p{0.22\linewidth}}
\toprule
ID & Scientific question & Main observable & Essential control \\
\midrule
\endhead
Q1 & Can multidimensional spectroscopy probe the spin sector through SOC? & Variation of \(S(\omega_1,\omega_3)\), cross peaks, integrated intensities, phases. & Turn off the spin--optical link: \(\lambda_\delta=\lambda_C=g_Q=0\). \\
Q2 & Is bright--dark coherence an indicator of quantum correlations? & Amplitude and phase of \(B\leftrightarrow D\) cross peaks, oscillations in \(t_2\). & Compare a model with fixed \(C_1\) to a model where \(C_1(T,B)\) varies. \\
Q3 & Are cross peaks a signature of coherent eigenstate mixing? & Scaling of cross peaks with \(V_{\mathrm{BD}}\), eigenvector analysis. & \(\mu_D=0\), \(V_{\mathrm{BD}}=0\): the cross peak must disappear. \\
Q4 & Which microscopic quantity does the cross-peak amplitude measure? & Derivatives \(\partial I/\partial \delta\), \(\partial I/\partial C_1\), \(\partial I/\partial g_Q\). & Independent scans of \(\delta\), \(C_1\), \(g_Q\). \\
Q5 & Does the magnetic field tune coherent mixing or only shift energies? & Peak shifts versus intensity changes. & Compare field dependence in \(H_{\mathrm{spin}}\), in \(T_{\mathrm{SP}}(B)\), and in the optical energies separately. \\
Q6 & Can nonlinear spectroscopy distinguish thermal and nonthermal states? & Spectra computed from thermal populations versus imposed populations. & Same Hamiltonian; only the initial distribution changes. \\
Q7 & Do cross peaks encode short-range spin correlations? & Correlation between \(I_{\mathrm{cross}}\) and \(C_1\). & Keep \(\delta\) fixed and vary only \(C_1\), then do the inverse. \\
Q8 & What is the role of detection in accessing the dark state? & Comparison of excitation/detection through \(\mu_B\), \(\mu_D\), and eigenstate mixing. & Compare \(\mu_D=0\) and \(\mu_D\neq0\). \\
Q9 & Are the observed correlations robust against numerical artifacts? & Convergence in \(N_k\), \(\eta\), frequency grid, normalization, chain size. & Reproduce the trends at several resolutions. \\
\bottomrule
\end{longtable}

\section{General validation principles}

\subsection{Add ingredients one by one}

The recommended validation order is:
\begin{equation}
\begin{aligned}
&\text{minimal bright/dark}
\rightarrow
\text{bosonic mode}
\rightarrow
\text{spin observables}
\\
&\rightarrow
k\text{-dispersion}
\rightarrow
\text{field/temperature}
\rightarrow
\text{explicit SOC}.
\end{aligned}
\end{equation}

Each step must keep the previous-step results as controls. A new signature is interpretable only if it is absent in the corresponding control.

\subsection{Always keep non-normalized spectra}

Panel-normalized figures are useful for comparing shapes. They do not allow conclusions about absolute intensity. For each spectrum, one should save at least
\begin{equation}
S_{R1}^{\mathrm{real}},\quad
S_{R1}^{\mathrm{imag}},\quad
|S_{R1}|,\quad
S_{R2}^{\mathrm{real}},\quad
S_{R2}^{\mathrm{imag}},\quad
|S_{R2}|,
\end{equation}
with and without normalization.

\subsection{Measure scalars, not only contour plots}

The minimal observables are:
\begin{align}
I_{\max} &= \max_{\omega_1,\omega_3}|S(\omega_1,\omega_3)|,\\
I_{\mathrm{int}} &= \int d\omega_1\,d\omega_3\,|S(\omega_1,\omega_3)|,\\
I_{\mathrm{cross}} &= \int_{\Omega_{\mathrm{cross}}} d\omega_1\,d\omega_3\,|S(\omega_1,\omega_3)|,\\
I_{\mathrm{diag}} &= \int_{\Omega_{\mathrm{diag}}} d\omega_1\,d\omega_3\,|S(\omega_1,\omega_3)|,\\
R_{\mathrm{cross}} &= I_{\mathrm{cross}}/I_{\mathrm{diag}}.
\end{align}

The peak positions should also be saved:
\begin{equation}
(\omega_1^\ast,\omega_3^\ast)
=
\arg\max_{\omega_1,\omega_3}|S(\omega_1,\omega_3)|.
\end{equation}

\section{Step 0 -- Numerical hygiene and conventions}

\subsection{Objective}

Before interpreting the physics, one must verify that the numerical conventions do not create the signatures. This step does not yet answer the physical question; it prevents building an interpretation on an artifact.

\subsection{Tests}

\testbox{\textbf{T0.1 -- Normalization control.} For the same parameter set, produce spectra with global normalization, panel normalization, and no normalization. The shapes may look similar, but conclusions about intensity must be drawn only from non-normalized spectra.}

\testbox{\textbf{T0.2 -- Frequency-grid convergence.} Recompute a representative spectrum with at least three grid resolutions. Peak positions and integrals over fixed windows must converge.}

\testbox{\textbf{T0.3 -- Broadening convergence in \(\eta\).} Test several values of \(\eta\). A robust physical signature should deform continuously with \(\eta\), not appear or disappear discontinuously.}

\testbox{\textbf{T0.4 -- \(R1\) and \(R2\) pathway control.} Save \(R1\), \(R2\), and their sum separately. A variation assigned to physics must be localizable: either it affects both pathways coherently, or it reveals an interpretable asymmetry.}

\decisionbox{Proceed to Step 1 only if the non-normalized spectra, peak positions, and integrated intensities are numerically stable for the minimal model.}

\section{Step 1 -- Minimal optical system without mode and without spin}

\subsection{Minimal system}

The first system should be as small as possible:
\begin{equation}
\mathcal{H}_0
=
\mathrm{span}\{\ket{g},\ket{B},\ket{D}\}.
\end{equation}
Impose:
\begin{equation}
N_k=1,\qquad
g_Q=0,\qquad
\omega_Q \text{ absent},\qquad
\mu_D=0.
\end{equation}
The optical Hamiltonian is
\begin{equation}
H_0
=
E_B\op{B}{B}
+
E_D\op{D}{D}
+
V_{\mathrm{BD}}L_{\mathrm{BD}}.
\end{equation}

\subsection{Why start here}

This system isolates the most fundamental mechanism: does a dark state become visible because it is mixed with a bright state? If this step is not understood, scans in \(\delta\), \(C_1\), \(k\), and \(g_Q(k)\) will be ambiguous.

\subsection{Tests}

\testbox{\textbf{T1.1 -- Strictly invisible dark state.} Set \(\mu_D=0\) and \(V_{\mathrm{BD}}=0\). The spectrum should contain only signatures associated with \(\ket{B}\). Any signature of \(\ket{D}\) in this control indicates a problem with the dipole, the basis, the spectral window, or the optical pathway.}

\testbox{\textbf{T1.2 -- Controlled appearance of the dark peak.} Scan \(V_{\mathrm{BD}}\) on a logarithmic scale, for example from \(10^{-5}\) to \(10^{-1}\) eV if numerically reasonable. Measure \(I_D\), \(I_{\mathrm{cross}}\), and \(R_{\mathrm{cross}}\).}

\testbox{\textbf{T1.3 -- Mixing scaling law.} For \(|V_{\mathrm{BD}}|\ll|E_D-E_B|\), the bright weight of the dark state should have the order of magnitude
\begin{equation}
\left|\frac{V_{\mathrm{BD}}}{E_D-E_B}\right|^2.
\end{equation}
The log--log slope of the dark intensity or cross-peak intensity should therefore be measured empirically. Depending on the optical pathway and the chosen observable, the slope can be linear or quadratic in \(V_{\mathrm{BD}}\); it should be documented rather than assumed.}

\testbox{\textbf{T1.4 -- Detuning dependence.} Repeat the scan for several \(\Delta_{BD}=E_D-E_B\). If the dark signal truly comes from mixing, it should increase when \(|\Delta_{BD}|\) decreases, all else being equal.}

\testbox{\textbf{T1.5 -- Sign of \(V_{\mathrm{BD}}\).} Compare \(+V_{\mathrm{BD}}\) and \(-V_{\mathrm{BD}}\). Absolute intensities may be nearly identical, but real/imaginary parts and phases can change. This comparison helps distinguish an intensity signature from a coherence signature.}

\decisionbox{This step partially answers Q3. If the cross peaks disappear when \(V_{\mathrm{BD}}=0\) and grow systematically with \(V_{\mathrm{BD}}\), then they are compatible with coherent bright--dark mixing. Otherwise, they cannot be interpreted as evidence of SOC or spin correlations.}

\section{Step 2 -- Add a local mode without \texorpdfstring{\(k\)}{k} dependence}

\subsection{System}

Add a truncated mode:
\begin{equation}
\mathcal{H}_1
=
\mathrm{span}\{\ket{g},\ket{B},\ket{D}\}
\otimes
\mathcal{H}_{\mathrm{ph}},
\end{equation}
but keep:
\begin{equation}
N_k=1,\qquad
\omega_Q(k)=\omega_Q,\qquad
g_Q(k)=g_Q.
\end{equation}

The Hamiltonian becomes:
\begin{equation}
H_1
=
E_B\op{B}{B}
+
E_D\op{D}{D}
+
\omega_Q a^\dagger a
+
V_{\mathrm{BD}}L_{\mathrm{BD}}
+
g_Q L_{\mathrm{BD}}(a+a^\dagger).
\end{equation}

\subsection{Objective}

This step separates static mixing \(V_{\mathrm{BD}}\) from mode-assisted dynamic mixing \(g_Q\). It indicates whether spectral structures come from simple bright--dark hybridization or from coupling to a collective coordinate.

\subsection{Tests}

\testbox{\textbf{T2.1 -- Control without dynamic mode.} Set \(g_Q=0\) and compare to Step 1. Differences should come only from the larger Hilbert space, not from a new physical coupling.}

\testbox{\textbf{T2.2 -- Activation of \(g_Q\).} Scan \(g_Q\) at fixed \(V_{\mathrm{BD}}\). Measure the appearance of satellites, peak deformation, and phase changes in \(R1\) and \(R2\).}

\testbox{\textbf{T2.3 -- Control with \(V_{\mathrm{BD}}=0\), \(g_Q\neq0\).} This test is important: if the dynamic mode alone can produce a dark signal, then dark-state access does not depend only on static mixing. If the signal disappears, the mode mainly acts as a modulation of already-present mixing.}

\testbox{\textbf{T2.4 -- Bosonic truncation convergence.} Repeat the test for \(n_{\mathrm{bosons}}=2,3,4,\ldots\). Intensities and peak positions should converge in the energy window used.}

\decisionbox{This step answers the difference between static mixing and dynamic coupling. One should not introduce \(k\) before identifying what \(g_Q\) already does in a single-mode model.}

\section{Step 3 -- Connect \texorpdfstring{\(\delta\)}{delta} and \texorpdfstring{\(C_1\)}{C1} without dispersion}

\subsection{System}

Keep \(N_k=1\) and replace the free parameter \(V_{\mathrm{BD}}\) by
\begin{equation}
V_{\mathrm{BD}}
=
V_0+\lambda_\delta\delta+\lambda_C C_1.
\end{equation}
At first, \(\delta\) and \(C_1\) should be controlled manually, without yet using the full spin chain.

\subsection{Objective}

This step directly addresses Q4 and Q7: does the spectrum measure \(\delta\), \(C_1\), or only the effective combination \(V_{\mathrm{BD}}\)?

\subsection{Tests}

\testbox{\textbf{T3.1 -- Scan \(\delta\) only.} Set \(C_1=C_1^{\mathrm{ref}}\), \(\lambda_C=0\), and scan \(\delta\). Measure \(I_{\mathrm{cross}}(\delta)\), \(I_D(\delta)\), and peak positions.}

\testbox{\textbf{T3.2 -- Scan \(C_1\) only.} Set \(\delta=\delta^{\mathrm{ref}}\), \(\lambda_\delta=0\), and scan \(C_1\). Measure the same observables.}

\testbox{\textbf{T3.3 -- Parameter degeneracy.} Compare two pairs \((\delta,C_1)\) that give the same \(V_{\mathrm{BD}}\). If the spectra are identical, the optical model does not distinguish \(\delta\) and \(C_1\); it measures only \(V_{\mathrm{BD}}\).}

\testbox{\textbf{T3.4 -- Local sensitivities.} Compute numerically:
\begin{equation}
\frac{\partial I_{\mathrm{cross}}}{\partial \delta},
\qquad
\frac{\partial I_{\mathrm{cross}}}{\partial C_1},
\qquad
\frac{\partial \phi_{\mathrm{cross}}}{\partial \delta},
\qquad
\frac{\partial \phi_{\mathrm{cross}}}{\partial C_1}.
\end{equation}
Phases may contain information different from intensity.}

\decisionbox{If two pairs \((\delta,C_1)\) with the same \(V_{\mathrm{BD}}\) produce the same spectrum, then spectroscopy does not separately measure \(\delta\) and \(C_1\) in this model. To separate them, one will need to add a channel where \(\delta\) and \(C_1\) affect different terms, for example \(\omega_Q(k)\), \(g_Q(k)\), or dephasing rates.}

\section{Step 4 -- Introduce the finite spin chain}

\subsection{System}

Now compute \(\delta(T,B)\), \(C_1(T,B)\), and \(\Delta_{\mathrm{spin}}\) from the spin sector:
\begin{equation}
H_{\mathrm{spin}}
=
\sum_i J_i\bm{S}_i\cdot\bm{S}_{i+1}
+
J_2\sum_i\bm{S}_i\cdot\bm{S}_{i+2}
+
g\mu_BB\sum_iS_i^z.
\end{equation}

\subsection{Objective}

The objective is to verify that the temperature/field dependence of the spectrum truly comes from spin observables, not from an arbitrary variation of optical parameters.

\subsection{Tests}

\testbox{\textbf{T4.1 -- Validation of \(C_1\).} Before computing any optical spectrum, plot \(C_1(T)\), \(C_1(B)\), \(\delta(T,B)\), and \(\Delta_{\mathrm{spin}}(T,B)\). Verify that they vary in the physically expected direction.}

\testbox{\textbf{T4.2 -- Chain size.} Repeat the spin observables for several \(N_{\mathrm{spin}}\). The qualitative trends of \(C_1\) and \(\Delta_{\mathrm{spin}}\) should be robust.}

\testbox{\textbf{T4.3 -- Boundary conditions.} Compare open and periodic boundary conditions if the code allows it. If \(C_1\) changes strongly, the optical interpretation must include finite-size uncertainty.}

\testbox{\textbf{T4.4 -- Optical control at fixed \(V_{\mathrm{BD}}\).} Vary \(T\) and \(B\) in the spin chain, but force \(V_{\mathrm{BD}}\) to remain constant in the optical model. If the spectrum still changes, there is another \(T,B\)-dependent channel that must be identified.}

\testbox{\textbf{T4.5 -- Active spin--optical control.} Compare with the full case where \(V_{\mathrm{BD}}(T,B)=V_0+\lambda_\delta\delta(T,B)+\lambda_CC_1(T,B)\). The difference between T4.4 and T4.5 isolates the spin--optical coupling effect.}

\decisionbox{This step is the first real test of Q1 and Q7. If the spectral changes follow \(C_1(T,B)\) or \(\delta(T,B)\) monotonically and disappear when \(\lambda_\delta=\lambda_C=0\), then the model supports the idea that the optical response probes the spin sector.}

\section{Step 5 -- Magnetic field: coherent mixing or simple energy shift}

\subsection{Objective}

Question Q5 asks whether the magnetic field tunes coherent mixing or only shifts energies. We must therefore artificially decompose the channels through which \(B\) enters the model.

\subsection{Scenarios}

\begin{longtable}{>{\centering\arraybackslash}p{0.12\linewidth} p{0.36\linewidth} p{0.38\linewidth}}
\toprule
Scenario & What depends on \(B\) & Interpretation tested \\
\midrule
\endhead
B0 & Nothing. & Numerical control. \\
B1 & \(T_{\mathrm{SP}}(B)\) and \(\delta(T,B)\) only. & The field affects dimerization. \\
B2 & \(C_1(B)\) only. & The field affects magnetic correlations. \\
B3 & Optical energies \(E_B,E_D\) only. & The field shifts resonances without changing mixing. \\
B4 & Everything active. & Full model. \\
\bottomrule
\end{longtable}

\subsection{Tests}

\testbox{\textbf{T5.1 -- Shift versus intensity.} For each scenario, extract \((\omega_1^\ast,\omega_3^\ast)\), \(I_{\mathrm{cross}}\), \(I_{\mathrm{diag}}\), and \(R_{\mathrm{cross}}\). A pure energy shift should mainly change positions; a change in mixing should mainly change relative intensities and phases.}

\testbox{\textbf{T5.2 -- Cross-peak phase.} The field may modify the phase even if the intensity changes little. Real and imaginary parts must therefore be retained, not only \(|S|\).}

\testbox{\textbf{T5.3 -- Field derivatives.} Compute
\begin{equation}
\frac{dI_{\mathrm{cross}}}{dB},
\qquad
\frac{d\omega^\ast}{dB},
\qquad
\frac{dR_{\mathrm{cross}}}{dB}.
\end{equation}
These three derivatives distinguish coupling modulation from a simple spectral shift.}

\decisionbox{If \(B\) mainly changes peak positions, the model indicates an energetic effect. If \(B\) mainly changes \(R_{\mathrm{cross}}\), phases, or the \(R1/R2\) ratio, there is evidence for modulation of coherent mixing.}

\section{Step 6 -- Temperature and Spin-Peierls transition}

\subsection{Objective}

The Spin-Peierls transition provides a natural control: above \(T_{\mathrm{SP}}\), \(\delta=0\). If the signal attributed to dimerization persists unchanged above \(T_{\mathrm{SP}}\), the attribution is weak.

\subsection{Tests}

\testbox{\textbf{T6.1 -- Scan around \(T_{\mathrm{SP}}\).} Compute spectra for \(T<T_{\mathrm{SP}}\), \(T\simeq T_{\mathrm{SP}}\), and \(T>T_{\mathrm{SP}}\). Extract \(I_{\mathrm{cross}}\), \(R_{\mathrm{cross}}\), and phases.}

\testbox{\textbf{T6.2 -- Separating \(\delta\) from \(C_1\).} Above \(T_{\mathrm{SP}}\), \(\delta=0\) but \(C_1\) may remain nonzero. This region tests whether the spectrum is sensitive to spin correlations without dimerization.}

\testbox{\textbf{T6.3 -- Control with \(\lambda_C=0\).} If \(\lambda_C=0\), the spin-dependent signal should collapse with \(\delta\) at the transition. If a signal persists, it comes from another mechanism.}

\testbox{\textbf{T6.4 -- Control with \(\lambda_\delta=0\).} If \(\lambda_\delta=0\), the spectrum can remain sensitive to \(C_1\) even above \(T_{\mathrm{SP}}\). This is the central test for whether short-range correlations are optically visible.}

\decisionbox{This step answers Q2 and Q7. The strongest result would be a signature that clearly distinguishes long-range dimerization \(\delta\) from short-range correlations \(C_1\).}

\section{Step 7 -- Add spin-sector dispersion}

\subsection{System}

Activate:
\begin{equation}
N_k>1,
\qquad
s=\texttt{spin\_mode\_dispersion\_scale}>0,
\qquad
A_Q=0.
\end{equation}
The energies \(E_B\) and \(E_D\) should initially remain independent of \(k\).

\subsection{Objective}

This step tests the idea that \(k\) mainly represents the spin/Spin-Peierls sector, not an optical \(d\)-band. It should show whether the dispersive spin dynamics leaves a measurable imprint on the 2D spectrum.

\subsection{Tests}

\testbox{\textbf{T7.1 -- Convergence in \(N_k\).} Compute \(N_k=1,5,10,25,50,100\). Integrated intensities must converge. If the structure changes strongly between \(50\) and \(100\), the model is not yet numerically stable.}

\testbox{\textbf{T7.2 -- Scan of \(s\).} Scan \(s=0,0.25,0.5,1,2\). The case \(s=0\) is the no-dispersion control.}

\testbox{\textbf{T7.3 -- \(k\)-resolved decomposition.} Save or reconstruct \(S_k(\omega_1,\omega_3)\) for several representative \(k\) values. This shows whether the integrated spectrum hides opposing contributions.}

\testbox{\textbf{T7.4 -- Keep optical energies flat.} Keep \(t_{B,r}=t_{D,r}=0\). If optical dispersions are activated too early, one will no longer know whether changes come from spin or optical bands.}

\decisionbox{Spin-sector dynamics can be called visible only if \(s\) modifies scalar observables or spectral shapes robustly, and if this effect disappears when \(s=0\).}

\section{Step 8 -- Modulation of \texorpdfstring{\(g_Q(k)\)}{gQ(k)}}

\subsection{Objective}

The parameter \(A_Q=\texttt{g\_Q\_k\_modulation}\) does not necessarily move the peaks. It redistributes the dynamic coupling over the \(k\) sector. One should therefore look for variations in intensity, spectral weight, and \(k\)-resolved contributions, not only for visual changes in normalized contours.

\subsection{Tests}

\testbox{\textbf{T8.1 -- Non-normalized scan of \(A_Q\).} Use \(A_Q=0,0.001,0.005,0.01,0.05,0.1,0.25,0.5,0.75\). For each point, measure \(I_{\max}\), \(I_{\mathrm{int}}\), \(I_{\mathrm{cross}}\), \(R_{\mathrm{cross}}\), and \(I_{R1}/I_{R2}\).}

\testbox{\textbf{T8.2 -- Difference maps.} Compute
\begin{equation}
\Delta S(A_Q)=S(A_Q)-S(0).
\end{equation}
Difference maps must be non-normalized or normalized by a common scale.}

\testbox{\textbf{T8.3 -- Contributions by \(k\) region.} Compare contributions near \(k=0\), \(k=\pi/2\), and \(k=\pi\). Since \(g_Q(k)=g_Q(1+A_Q\cos k)\), \(A_Q\) strengthens one region and weakens the other.}

\testbox{\textbf{T8.4 -- Sign control.} Also test \(A_Q<0\). If the model is consistent, \(A_Q>0\) and \(A_Q<0\) should redistribute weight toward opposite \(k\) regions.}

\decisionbox{If \(A_Q\) progressively decreases the total intensity without moving the peaks, this is not a failure: it is compatible with spectral-weight modulation. To interpret it physically, one must show which \(k\) region loses or gains weight.}

\section{Step 9 -- Direct access to the dark state versus bright--dark mixing}

\subsection{Objective}

The discussion of \(d\)-\(d\) transitions, selection rules, and the crystal field shows that two mechanisms must be distinguished:
\begin{align}
\text{direct channel:}\quad & \mu_D\neq0,\quad V_0=0,\\
\text{mixing channel:}\quad & \mu_D=0,\quad V_0\neq0.
\end{align}

\subsection{Tests}

\testbox{\textbf{T9.1 -- Scenario A: direct dark dipole.} Set \(V_0=0\), \(\lambda_\delta=\lambda_C=0\), but \(\mu_D\neq0\). Measure dark and cross peaks.}

\testbox{\textbf{T9.2 -- Scenario B: bright--dark mixing.} Set \(\mu_D=0\), \(V_0\neq0\). Adjust \(V_0\) to obtain a dark intensity comparable to Scenario A.}

\testbox{\textbf{T9.3 -- Discrimination by pathways.} Compare \(R1\), \(R2\), phases, and peak ratios. Two mechanisms that give the same linear intensity may give different 2D responses.}

\testbox{\textbf{T9.4 -- Temperature/field dependence.} If the direct \(\mu_D\) channel is purely crystal-field-like, it should not strongly follow \(\delta(T,B)\) or \(C_1(T,B)\), unless an additional hypothesis is introduced. The spin-dependent mixing channel should follow these variables.}

\decisionbox{This step answers Q8 and clarifies the physical meaning of \(V_0\). If Scenarios A and B are indistinguishable in the chosen observables, the model cannot conclude that the dark signal comes from SOC rather than from a weakly allowed direct dipole.}

\section{Step 10 -- Thermal versus nonthermal states}

\subsection{Objective}

Question Q6 requires changing the initial population without changing the Hamiltonian. Otherwise, one does not know whether the difference comes from a nonthermal state or from a parameter change.

\subsection{Tests}

\testbox{\textbf{T10.1 -- Thermal reference state.} Use Boltzmann populations for a temperature \(T\).}

\testbox{\textbf{T10.2 -- Simple nonthermal population.} Build a distribution in which selected bright/dark levels or selected mode levels are overpopulated. Keep the same Hamiltonian and the same dipoles.}

\testbox{\textbf{T10.3 -- Comparison at equal average energy.} Compare a thermal state and a nonthermal state with similar average energy. If the spectra differ, spectroscopy is sensitive to the population structure, not only to energy.}

\testbox{\textbf{T10.4 -- \(t_2\) dependence.} If the code allows a population-time scan \(t_2\), look for oscillations or dephasings that differ between the thermal and nonthermal states.}

\decisionbox{Sensitivity to nonthermal states can be claimed only if two states with identical Hamiltonian parameters produce distinct and robust 2D signatures.}

\section{Step 11 -- Robustness and artifacts}

\subsection{Objective}

Before drawing conclusions, one must show that the trends do not depend on an arbitrary numerical choice.

\subsection{Tests}

\testbox{\textbf{T11.1 -- Robustness to \(\eta\).} Trends in \(I_{\mathrm{cross}}\), \(R_{\mathrm{cross}}\), and peak positions must survive reasonable variations of the broadening.}

\testbox{\textbf{T11.2 -- Grid robustness.} Repeat the main results with a finer frequency grid and a slightly wider spectral window.}

\testbox{\textbf{T11.3 -- Robustness to \(N_k\).} Conclusions about \(A_Q\), \(s\), and \(\omega_Q(k)\) must remain stable when \(N_k\) is increased.}

\testbox{\textbf{T11.4 -- Robustness to chain size.} Conclusions related to \(C_1\) and \(\Delta_{\mathrm{spin}}\) must be compared for several chain sizes.}

\testbox{\textbf{T11.5 -- Robustness to integration windows.} The integrals \(I_{\mathrm{cross}}\) and \(I_{\mathrm{diag}}\) must be recomputed with slightly different windows.}

\decisionbox{A scientific conclusion should be classified as robust only if it survives at least three controls: no panel normalization, numerical convergence, and a physical control in which the mechanism is switched off.}

\section{Step 12 -- Extension toward explicit SOC}

\subsection{Why not start there}

An explicit Hamiltonian
\begin{equation}
H_{\mathrm{SOC}}=\lambda\bm{L}\cdot\bm{S}
\end{equation}
is more physical, but it increases the Hilbert-space size and introduces new choices: \(L_\alpha\) matrices, orbital basis, dipoles, crystal field, and possible double counting with \(V_{\mathrm{BD}}\).

\subsection{Route 1: effective spinor model}

The first reasonable extension is:
\begin{equation}
(g,D,B)\otimes(\uparrow,\downarrow)\otimes\mathcal{H}_{\mathrm{ph}}.
\end{equation}
For \(n_{\mathrm{bosons}}=3\), the Hilbert dimension per \(k\) increases from \(9\) to \(18\), and the Liouville dimension from \(81\) to \(324\).

\testbox{\textbf{T12.1 -- Recover the effective model.} Choose simple \(L_\alpha^{\mathrm{eff}}\) matrices and verify that, at weak \(\lambda_{\mathrm{eff}}\), the signatures resemble the model with \(V_{\mathrm{BD}}\).}

\testbox{\textbf{T12.2 -- Avoid double counting.} When \(H_{\mathrm{SOC}}\) is active, replace \(V_{\mathrm{BD}}\) by a residual \(V_{\mathrm{res}}\), or set it to zero in the first test.}

\testbox{\textbf{T12.3 -- Spin-flip signatures.} Verify whether new pathways appear when \(S_x\) and \(S_y\) are active. These pathways would be absent from a scalar bright--dark model.}

\subsection{Route 2: microscopic \texorpdfstring{\(d\)}{d}-orbital model}

The more complete route is:
\begin{equation}
\{d_{x^2-y^2},d_{xy},d_{xz},d_{yz},d_{z^2}\}
\otimes
(\uparrow,\downarrow)
\otimes
\mathcal{H}_{\mathrm{ph}}.
\end{equation}
For \(n_{\mathrm{bosons}}=3\), the Hilbert dimension per \(k\) is about \(30\), and the Liouville dimension about \(900\).

\decisionbox{Route 2 should be undertaken only after identifying, with Route 1, which 2D observable is truly sensitive to explicit spin. Otherwise, the microscopic model risks adding parameters without increasing the strength of the conclusion.}

\section{Recommended calculation order}

\begin{longtable}{>{\centering\arraybackslash}p{0.08\linewidth} p{0.28\linewidth} p{0.32\linewidth} p{0.20\linewidth}}
\toprule
Priority & Calculation & Why & Target question \\
\midrule
\endhead
1 & \(V_{\mathrm{BD}}\) scan, \(\mu_D=0\), \(g_Q=0\), \(N_k=1\). & Prove that cross peaks come from mixing. & Q3 \\
2 & \(g_Q\) scan, \(N_k=1\). & Separate static mixing from dynamic mode coupling. & Q1, Q3 \\
3 & Manual scans of \(\delta\), \(C_1\). & Determine whether the spectrum measures only \(V_{\mathrm{BD}}\) or distinguishes its components. & Q4, Q7 \\
4 & Finite spin chain: \(C_1(T,B)\), \(\delta(T,B)\). & Connect the optical model to spin observables. & Q1, Q7 \\
5 & Temperature scan around \(T_{\mathrm{SP}}\). & Use the transition as a physical control. & Q2, Q7 \\
6 & Magnetic-field decomposition. & Distinguish energetic shift from mixing change. & Q5 \\
7 & Activate \(\omega_Q(k)\), \(A_Q=0\). & Test dispersive spin dynamics. & Q1 \\
8 & Scan \(A_Q\). & Test redistribution of coupling in \(k\). & Q1, Q9 \\
9 & Compare \(\mu_D\neq0\) versus \(V_0\neq0\). & Distinguish direct dipole from hybridization. & Q8 \\
10 & Nonthermal states. & Test sensitivity to populations. & Q6 \\
11 & Effective spinor model. & Replace \(V_{\mathrm{BD}}\) by explicit SOC. & Q1--Q8 \\
\bottomrule
\end{longtable}

\section{Conclusion criteria}

\subsection{Strong conclusion}

One can strongly support the claim that multidimensional spectroscopy probes the spin sector if the following conditions are satisfied:
\begin{enumerate}
    \item Cross peaks disappear when bright--dark mixing and the direct dark channel are both switched off.
    \item Cross peaks systematically follow \(\delta\), \(C_1\), or \(\omega_Q(k)\) when these quantities are activated separately.
    \item The trends persist without panel normalization.
    \item The trends converge in grid, \(N_k\), \(\eta\), and chain size.
    \item A control with \(\lambda_\delta=\lambda_C=g_Q=0\) removes the spin-dependent signature.
\end{enumerate}

\subsection{Moderate conclusion}

The conclusion will be moderate if the spectra are sensitive to \(V_{\mathrm{BD}}\), but do not separate \(\delta\) and \(C_1\). In that case, the model shows that spectroscopy can probe an effective mixing parameter, but cannot yet identify its microscopic source.

\subsection{Weak or negative conclusion}

The conclusion will be weak if:
\begin{enumerate}
    \item the same signatures appear with \(\mu_D\neq0\) and without spin coupling;
    \item the signatures disappear under small changes of normalization or grid;
    \item field or temperature variations only shift peaks without changing cross-peak ratios;
    \item effects assigned to \(k\) do not converge with \(N_k\).
\end{enumerate}

\section{Recommended deliverables}

For each important step, save:
\begin{enumerate}
    \item raw spectra in numerical format;
    \item non-normalized figures;
    \item normalized figures only for visual comparison;
    \item a complete parameter file;
    \item a CSV table of scalar observables;
    \item a short note indicating which mechanism was active or inactive.
\end{enumerate}

The minimal observable table should contain:
\begin{equation}
I_{\max},\quad
I_{\mathrm{int}},\quad
I_{\mathrm{cross}},\quad
I_{\mathrm{diag}},\quad
R_{\mathrm{cross}},\quad
I_{R1},\quad
I_{R2},\quad
\omega_1^\ast,\quad
\omega_3^\ast.
\end{equation}

\section{Next concrete action}

The next calculation should be Step 1:
\begin{equation}
N_k=1,\qquad
g_Q=0,\qquad
\mu_D=0,\qquad
V_{\mathrm{BD}}\ \text{manually scanned}.
\end{equation}

This calculation is the smallest possible test of the core of the model. It answers the question:
\begin{quote}
Are the cross peaks that we observe truly generated by coherent bright--dark mixing?
\end{quote}

Without this validation, it would be premature to interpret scans in \(g_Q(k)\), \(\delta\), \(C_1\), \(T\), or \(B\) as signatures of the spin sector.

\end{document}
